\documentclass{article}
%  ╔══════════════════════════════════════════════════════════╗
%  ║                         Title                            ║
%  ╚══════════════════════════════════════════════════════════╝
%NOTE: I want \title,\author to be near top, and this \renewcommand{\maketitle} fails if these are above it.
\usepackage[utf8]{inputenc}         % Used to compile any UTF-8 Character (and supports Unicode)
\usepackage{titling}

\renewcommand{\maketitle}{
    \quad
    \vspace{1em}
  \begin{center}
      \LARGE\bfseries \thetitle \par
      \vspace{0.6em}
      \large
  \end{center}
    \vspace{1em}
}

\title{Langlands for SL2}
\author{Nathanael Chwojko-Srawley}
\date{\today}

%  ╔══════════════════════════════════════════════════════════╗
%  ║                         packages                         ║
%  ╚══════════════════════════════════════════════════════════╝

\usepackage{extarrows}         % Used to compile any UTF-8 Character (and supports Unicode)
\usepackage{stackengine}
\usepackage{colortbl}
\usepackage{scalerel}
\usepackage{amsmath, amssymb} % Used to write better mathematical expressions (ex. \mathbb{R})
\usepackage{stackrel}         % for left-right arrows, staking symbol above and bellow
\usepackage{mathrsfs}
\usepackage{bbm}                     % for mathbbm{1}
\usepackage{euscript}
\usepackage{commutative-diagrams}
\usepackage[font=itshape]{quoting}
\usepackage{mathtools}
  \usepackage{enumitem}
\usepackage{amsthm}           % Used to create theorem enviroments.
\usepackage{graphicx}         % Let's you use images in LaTeX; ex the \includegraphics command.
\usepackage{hhline}
%\usepackage{luatexja}
%\usepackage{fontspec}
\usepackage{dsfont}
\usepackage{wasysym}          % emoji's!
\usepackage{float}            % package with ``H'' for figures and tables, ex. Images, tables, etc
%\usepackage{subcaption}      % More options with sub-captions.
\usepackage{setspace}         % More flexibility on spacing in document.
\usepackage{hyperref}         % let's me put hyperlinks
%\usepackage{tikz}             % for commutative diagrams https://tex.stackexchange.com/questions/419221/how-to-do-the-pushout-with-universal-property
\usepackage{tikz-cd}          % for commutative diagrams https://tex.stackexchange.com/questions/419221/how-to-do-the-pushout-with-universal-property
%\usepackage{quiver}           % for better commutative diagrams
\usepackage{bookmark}         % creating heading shortcut bar on the left as in word
\usepackage{longtable}        % create tables that extend past one page
%\usepackage{siunitx}         % required for alignment
%\usepackage{multirow}        % for multiple rows
\usepackage{microtype}        % makes nice text. Fixes some under-/over- filled box problems.
%\usepackage{fancyvrb}        % Let's you write code (different style than listings).
\usepackage[most]{tcolorbox} 	      % Makes nice boxes for thms, cor, prop, or anything else.
%\usepackage{enumitem}        % More flexibility on enumeration (don't know much).
\usepackage{xcolor}           % Colour text.
\usepackage{wrapfig}          % To wrap text around figure
\usepackage{listings}         % Create nice colour code blocks (customized in preamble)
\usepackage{thmtools}         % Customize theorem environment (in preamble).
%\usepackage{marvosym}        % Provide ``basic'' symbols (hazard, religious, smiley face etc.)
\usepackage{array}            % \begin{table}{>{\em}c c} -> 1st column italic (bfseries for bold)
%\usepackage[english]{babel}  % If I want to blindtext in english.
\usepackage{blindtext}
\usepackage[a4paper, total={6in, 8.5in}]{geometry}
%\usepackage[symbol]{footmisc} % For daggers and asterisk for footnotes. 1: *, 2: dagger, etc.
\usepackage{fancyhdr}
\usepackage{titlesec}
%\usepackage{pst-plot}          %To make nice plots: https://tex.stackexchange.com/questions/47594/how-can-i-make-a-graph-of-a-function#47596
% \usepackage[answerdelayed]{exercise}
\usepackage{etoolbox}          % for Dynkin diagrams
\usepackage{dynkin-diagrams}   % for Dynkin diagrams

% ╔═════════════════════════════════════════════════════╗
% ║           for figures via inkspace                  ║
% ╚═════════════════════════════════════════════════════╝

% to import figures via: https://github.com/gillescastel/inkscape-figures
\usepackage{import}
\usepackage{pdfpages}
\usepackage{transparent}


% This command is the ONLY command imported via the packages snippet, think of it as a tiny package written out
\newcommand{\incfig}[2][1]{%
    \def\svgwidth{#1\columnwidth}
    \import{./figures/}{#2.pdf_tex}
}

% ╔═══════════════════════════════════════════════════════════╗
% ║                 bibliography and citing                   ║
% ╚═══════════════════════════════════════════════════════════╝

\usepackage{csquotes}
% \usepackage[backend=biber,citestyle=alphabetic,bibstyle=authortitle]{biblatex}
\usepackage{biblatex}

\addbibresource{bibliography.bib}

% ╔═══════════════════════════════════════════════════════════╗
% ║                    colours (colors)                       ║
% ╚═══════════════════════════════════════════════════════════╝

\definecolor{dkgreen}{rgb}{0,0.6,0}
\definecolor{gray}{rgb}{0.5,0.5,0.5}
\definecolor{mauve}{rgb}{0.58,0,0.82}
\definecolor{lightyellow}{rgb}{1,1,0.6}
\definecolor{reddish}{HTML}{A01A3A}

%  ╔══════════════════════════════════════════════════════════╗
%  ║                       book format                        ║
%  ╚══════════════════════════════════════════════════════════╝

\pagestyle{fancy}

\setlength\headheight{20pt}

\renewcommand{\sectionmark}[1]{\markright{\thesection.\ #1}}

\fancyfoot[C,CO]{\textcolor{gray}{Nathanael Chwojko-Srawley}}
\fancyfoot[R,RO]{\thepage}{}

% Create a new page style for the first page
\fancypagestyle{firstpage}{
  \fancyhf{} % Clear header and footer
  \fancyhead[L]{\theauthor}
  \fancyhead[R]{\today} % Date in the top-right
  \fancyfoot[C]{\textcolor{gray}{Nathanael Chwojko-Srawley}}
  \fancyfoot[R]{\thepage}
}

% Apply the first page style conditionally
\AtBeginDocument{\thispagestyle{firstpage}}

%have format the title command
\newcommand{\chapter}[1]{%
  \clearpage
  \vspace*{30pt} % Space before the chapter title
  {\normalfont\huge\bfseries #1\par} % Chapter title formatting
  \vspace{20pt} % Space after the chapter title
}


\setlength{\parindent}{0pt} % don't like indentation infront of paragraphs
\setlength{\parskip}{1ex}   %so that there is still space between paragarphs


%  ╔══════════════════════════════════════════════════════════╗
%  ║                       Shortcuts                          ║
%  ╚══════════════════════════════════════════════════════════╝

%\ExerciseLevelInToc

% I like original mathcal better, but need lowercase.
\DeclareFontFamily{OT1}{pzc}{}
\DeclareFontShape{OT1}{pzc}{m}{it}{<-> s * [1.10] pzcmi7t}{}
\DeclareMathAlphabet{\mathpzc}{OT1}{pzc}{m}{it}


% mathbb
\newcommand{\A}{{\mathbb{A}}}
\newcommand{\C}{{\mathbb{C}}}
\newcommand{\F}{{\mathbb{F}}}
\newcommand{\N}{{\mathbb{N}}}
\newcommand{\Q}{{\mathbb{Q}}}
\newcommand{\R}{{\mathbb{R}}}
\newcommand{\V}{{\mathbb{V}}}
\newcommand{\Z}{{\mathbb{Z}}}
\newcommand{\BI}{{\mathbb{I}}}
\renewcommand{\H}{{\mathbb{H}}}
\renewcommand{\P}{{\mathbb{P}}}


% mathcal
% \newcommand{\co}{{\mathfrak{o}}} %mathcal{o} looks like wreath product, so I changed it to mathfrak
\newcommand{\co}{{\mathpzc{o}}}
\newcommand{\CA}{{\mathcal{A}}}
\newcommand{\CB}{{\mathcal{B}}}
\newcommand{\CC}{{\mathcal{C}}}
\newcommand{\CD}{{\mathcal{D}}}
\newcommand{\CF}{{\mathcal{F}}}
\newcommand{\CG}{{\mathcal{G}}}
\newcommand{\CH}{{\mathcal{H}}}
\newcommand{\CI}{{\mathcal{I}}}
\newcommand{\CK}{{\mathcal{K}}}
\newcommand{\CL}{{\mathcal{L}}}
\newcommand{\CM}{{\mathcal{M}}}
\newcommand{\CN}{{\mathcal{N}}}
\newcommand{\CO}{{\mathcal{O}}}
\newcommand{\CQ}{{\mathcal{Q}}}
\newcommand{\CR}{{\mathcal{R}}}
\newcommand{\CS}{{\mathcal{S}}}
\newcommand{\CT}{{\mathcal{T}}}
\newcommand{\CV}{{\mathcal{V}}}
\newcommand{\CW}{{\mathcal{W}}}
\newcommand{\CZ}{{\mathcal{Z}}}
% EXCPETION:
\newcommand{\CP}{{\mathcal{P}}}
% \newcommand{\CP}{\mathbb{C}\mathrm{P}}
\newcommand{\RP}{\mathbb{R}\mathrm{P}}


%Euscript
\newcommand{\EB}{\EuScript{B}}
\newcommand{\EC}{{\EuScript{C}}}
\newcommand{\EF}{{\EuScript{F}}}
\newcommand{\EL}{{\EuScript{L}}}
\newcommand{\EO}{{\EuScript{O}}}


%mathfrak (lowercase)
\newcommand{\fa}{{\mathfrak{a}}}
\newcommand{\fb}{{\mathfrak{b}}}
\newcommand{\fc}{{\mathfrak{c}}}
\newcommand{\fd}{{\mathfrak{d}}}
\newcommand{\fm}{{\mathfrak{m}}}
\newcommand{\fn}{{\mathfrak{n}}}
\newcommand{\fo}{{\mathfrak{o}}}
\newcommand{\fp}{{\mathfrak{p}}}
\newcommand{\fq}{{\mathfrak{q}}}


%mathfrak (upper case)
\newcommand{\FA}{{\mathfrak{A}}}
\newcommand{\FB}{{\mathfrak{B}}}
\newcommand{\FC}{{\mathfrak{C}}}
\newcommand{\FD}{{\mathfrak{D}}}
\newcommand{\FK}{{\mathfrak{K}}}
\newcommand{\FM}{{\mathfrak{M}}}
\newcommand{\FN}{{\mathfrak{N}}}
\newcommand{\FO}{{\mathfrak{O}}}
\newcommand{\FP}{{\mathfrak{P}}}


%mathscr
\newcommand{\ff}{{\mathscr{F}}}
\newcommand{\SA}{\mathscr{A}}
\newcommand{\SC}{\mathscr{C}}
\newcommand{\SF}{\mathscr{F}}
\newcommand{\SG}{\mathscr{G}}
\newcommand{\SH}{\mathscr{H}}
\newcommand{\SI}{\mathscr{I}}
\newcommand{\SK}{\mathscr{K}}
\newcommand{\SN}{\mathscr{N}}
\newcommand{\SQ}{\mathscr{Q}}
\newcommand{\SR}{\mathscr{R}}
\newcommand{\SV}{\mathscr{V}}
\newcommand{\SZ}{\mathscr{Z}}


% operatorname -- 2 letters (lowercase and upper case)
\newcommand{\ab}{{\operatorname{ab}}}
\newcommand{\ad}{{\operatorname{ad}}}
\newcommand{\ch}{{\operatorname{ch}}}
\newcommand{\ev}{{\operatorname{ev}}}
\newcommand{\id}{{\operatorname{id}}}
\newcommand{\im}{{\operatorname{im}}}
\newcommand{\nm}{{\operatorname{Nm}}}
\newcommand{\op}{{\operatorname{op}}}
\newcommand{\rk}{{\operatorname{rk}}}
\newcommand{\sh}{{\operatorname{sh}}}
\newcommand{\tr}{{\operatorname{tr}}}

\newcommand{\Ab}{{\operatorname{Ab}}}
\newcommand{\Br}{{\operatorname{Br}}}
\newcommand{\Ch}{{\operatorname{Ch}}}
\newcommand{\Cl}{{\operatorname{Cl}}}
\newcommand{\GL}{{\operatorname{GL}}}
\newcommand{\Nm}{{\operatorname{Nm}}}
\newcommand{\SL}{{\operatorname{SL}}}

\renewcommand{\Re}{{\operatorname{Re}}}
\renewcommand{\Im}{{\operatorname{Im}}}


%categories
\newcommand{\Grp}{{\textbf{Grp}}}
\newcommand{\Fld}{{\textbf{Fld}}}
\newcommand{\Sh}{{\textbf{Sh}}}
\newcommand{\Set}{{\textbf{Set}}}
\newcommand{\Mod}{{\textbf{Mod}}}
\newcommand{\PSh}{{\textbf{PSh}}}
\newcommand{\Sch}{{\textbf{Sch}}}
\newcommand{\Top}{{\textbf{Top}}}
\newcommand{\RgSp}{{\textbf{RgSp}}}
\newcommand{\Ring}{{\textbf{Ring}}}


%operatorname -- 3+ letters (lowercase and upper case)
\newcommand{\rad}{{\operatorname{rad}}}
\newcommand{\sep}{{\operatorname{sep}}}
\newcommand{\res}{{\operatorname{res}}}
\newcommand{\sgn}{{\operatorname{sgn}}}
\newcommand{\pre}{{\operatorname{pre}}}
\newcommand{\ord}{{\operatorname{ord}}}
\newcommand{\vol}{{\operatorname{vol}}}
\newcommand{\lcm}{{\operatorname{lcm}}}
\newcommand{\len}{{\operatorname{len}}}

\newcommand{\conj}{{\operatorname{conj}}}
\newcommand{\disc}{{\operatorname{disc}}}
\newcommand{\stab}{{\operatorname{stab}}}
\newcommand{\kdim}{{\operatorname{kdim}}}
\newcommand{\diag}{{\operatorname{diag}}}

\newcommand{\trdeg}{\operatorname{trdeg}}
\newcommand{\coker}{{\operatorname{coker}}}
\newcommand{\codim}{{\operatorname{codim}}}

\newcommand{\Ann}{{\operatorname{Ann}}}
\newcommand{\Aut}{{\operatorname{Aut}}}
\newcommand{\Der}{{\operatorname{Der}}}
\newcommand{\Div}{{\operatorname{Div}}}
\newcommand{\Emb}{{\operatorname{Emb}}}
\newcommand{\End}{{\operatorname{End}}}
\newcommand{\Ext}{{\operatorname{Ext}}}
\newcommand{\Gal}{{\operatorname{Gal}}}
\newcommand{\Hom}{{\operatorname{Hom}}}
\newcommand{\Ind}{{\operatorname{Ind}}}
\newcommand{\Inn}{{\operatorname{Inn}}}
\newcommand{\Int}{{\operatorname{int}}}
\newcommand{\Irr}{{\operatorname{Irr}}}
\newcommand{\Jac}{{\operatorname{Jac}}}
\newcommand{\Mor}[1]{\operatorname{Mor}(\textbf{#1})}
\newcommand{\Nil}{{\operatorname{Nil}}}
\newcommand{\Obj}{{\operatorname{Obj}}}
\newcommand{\Out}{{\operatorname{Out}}}
\newcommand{\Pic}{{\operatorname{Pic}\,}}
\newcommand{\Rep}{{\operatorname{Rep}}}
\newcommand{\Res}{{\operatorname{Res}}}
\newcommand{\Spl}{{\operatorname{Spl}}}
\newcommand{\Syl}{{\operatorname{Syl}}}
\newcommand{\Sym}{{\operatorname{Sym}}}
\newcommand{\Tor}{{\operatorname{Tor}}}

\newcommand{\Char}{{\operatorname{char}}}
\newcommand{\Frac}{{\operatorname{Frac}}}
\newcommand{\Frob}{{\operatorname{Frob}}}
\newcommand{\Proj}{{\operatorname{Proj}\,}}
\newcommand{\RMod}{{}_R\operatorname{Mod}}
\newcommand{\SExt}{{\mathcal{E}\emph{xt}}}
\newcommand{\SHom}{{\emph{Hom}}}
\newcommand{\Span}{{\operatorname{span}}}
\newcommand{\Spec}{{\operatorname{Spec}\,}}
\newcommand{\Supp}{{\operatorname{Supp}}}
\newcommand{\Tens}[1]{#1\otimes --}
\newcommand{\fHom}[1]{\operatorname{Hom}(#1, --)}

\newcommand{\mSpec}{{\operatorname{mSpec}}}
\newcommand{\cofHom}[1]{\operatorname{Hom}(--, #1)}


% connectors
% \renewcommand{\mid}{\big|}
\newcommand{\sse}{\subseteq}
\newcommand{\subg}{\leqslant}
\newcommand{\supg}{\geqslant}
\newcommand{\ssne}{\subsetneq}
\newcommand{\isosubg}{\lesssim}
\newcommand{\nsse}{\not\subseteq}
\newcommand{\nsubg}{\trianglelefteq}
\newcommand{\nsupg}{\trianglerighteq}
\newcommand{\csubg}{\blacktriangleleft}
\renewcommand{\mid}{\big|}


% Arrows
\newcommand{\rw}{\rightarrow}
\newcommand{\Rw}{\Rightarrow}
\newcommand{\lw}{\leftarrow}
\newcommand{\Lw}{\Leftarrow}
\newcommand{\lrw}{\leftrightarrow}
\newcommand{\LRw}{\Leftrightarrow}
\newcommand{\acton}{\curvearrowright}
\newcommand{\actson}{\curvearrowright}
\newcommand\mapsfrom{\mathrel{\reflectbox{\ensuremath{\mapsto}}}}

% arrows for commutative diagram: create four new angled double-struck arrows
\newcommand\myrot[1]{\mathrel{\rotatebox[origin=c]{#1}{$\Rightarrow$}}}
\newcommand\NEarrow{\myrot{45}}
\newcommand\NWarrow{\myrot{135}}
\newcommand\SWarrow{\myrot{-135}}
\newcommand\SEarrow{\myrot{-45}}


%shorthands
\newcommand{\oln}[1]{{\overline{#1}}}
\renewcommand{\bf}[1]{\textbf{#1}}
\newcommand{\qn}{\quad\newline}


% limits
\newcommand{\Lim}{{\varprojlim}}
\newcommand{\coLim}{{\varinjlim}}
\newcommand{\Dlim}{\lim\limits_{\longrightarrow}}
\newcommand{\coDlim}{\lim\limits_{\longleftarrow}}


%for saturated ideals in the localization section (I don't like these, I think I'll delete)
\newcommand{\LIm}{{}^e }
\newcommand{\LPre}{{}^c }

%  ╔══════════════════════════════════════════════════════════╗
%  ║                    Custom Commands                       ║
%  ╚══════════════════════════════════════════════════════════╝

% swap varphi and phi
\let\temp\phi
\let\phi\varphi
\let\varphi\temp

\newcommand{\placeholder}{\_\_\_\_\_}

\newcommand{\set}[2]{\left\{#1 \ \middle|\ #2\right\}}

%mod should not have the crazy amount of space to its right!
\renewcommand{\mod}{\,\operatorname{mod}\,}

% dcups
\newcommand{\dcup}{\sqcup}
\newcommand{\Dcup}{\coprod}
% \newcommand{\Dcup}{\bigsqcup}

\newcommand{\nbot}{\mathrel{\rlap{$\bot$}\mkern2mu/}}

%a nice large circle
\newcommand{\tikzcircle}[2][red,fill=black]{\tikz[baseline=-0.5ex]\draw[#1,radius=#2] (0,0) circle ;}%

\makeatletter
\newcommand*\bigcdot{\mathpalette\bigcdot@{.5}}
\newcommand*\bigcdot@[2]{\mathbin{\vcenter{\hbox{\scalebox{#2}{$\m@th#1\bullet$}}}}}
\makeatother

%somehow not a default command
\def\rddots{\mathstrut^{.^{.^{.}}}}

%% new delimiter
\DeclarePairedDelimiter{\ceil}{\lceil}{\rceil}


%  ╔══════════════════════════════════════════════════════════╗
%  ║                     environments                         ║
%  ╚══════════════════════════════════════════════════════════╝


%remember the option: breakable

% create tcolor boxes
\newtcbtheorem[number within=section]{thm}{Theorem}%
{
colback=green!5,
colframe=green!35!black,
fonttitle=\bfseries,
label type=theorem
}{th}

\newtcbtheorem[number within=section, use counter from=thm]{lem}{Lemma}%
{
colback=blue!5,
colframe=black!35!black,
fonttitle=\bfseries,
label type=lemma
}{lm}
\newtcbtheorem[number within=section, use counter from=thm]{prop}{Proposition}%
{
colback=white!5,
colframe=white!35!black,
fonttitle=\bfseries,
label type=proposition
}{pr}
\newtcbtheorem[number within=section, use counter from=thm]{cor}{Corollary}%
{colback=blue!5,
colframe=blue!35!black,
fonttitle=\bfseries,
label type=corollary
}{co}
\newtcbtheorem[number within=section, use counter from=thm]{axiom}{Axiom}%
{colback=black!5,
colframe=yellow!35!black,
fonttitle=\bfseries,
label type=axiom
}{ax}
\newtcbtheorem[number within=section, use counter from=thm]{defn}{Definition}%
{colback=red!5,
colframe=red!35!black,
fonttitle=\bfseries,
label type=definition
}{df}


% for <s-cr> to create \item[$\LRw$] instead of \item
\newenvironment{equivEnumerate}
 {\begin{enumerate}
	}{
\end{enumerate}}



%Custom enviroments
 \newenvironment{note}
 {\begin{tcolorbox}[enhanced, sharp corners, colback=white , borderline={0.2pt}{0pt}{black}]
   \begin{quote}\textbf{Note}:
   }{
   \end{quote}\end{tcolorbox}}

 \newenvironment{intuition}
 {\begin{quote}\textbf{Intuition}:
   }{
   \end{quote}}

\newtcbtheorem[number within=section, use counter from=thm]{titledBox}{}{%
  enhanced,
  sharp corners,
  colback=white,
  colframe=black,
  borderline={0.2pt}{0pt}{black},
  left=2mm,
  right=2mm,
  top=2mm,
  bottom=2mm,
  title style={opacity=1},
  attach title to upper,
  before upper={\textbf{\tcbtitletext}\quad},
  coltitle=black,
  fonttitle=\bfseries,
  separator sign={\quad}
}{box}


%better example and proof environments
\def\exampletext{Example} % If English
\colorlet{colexam}{red!55!black} % Global example color


\newtcbtheorem[number within=section, use counter from=thm]{example}{Example}{% \newtcolorbox[use counter=testexample]{testexamplebox}{%
% Example Frame Start
empty,% Empty previously set parameters
title={\exampletext \thetcbcounter #1},% use \thetcbcounter to access the testexample counter text
% Attaching a box requires an overlay
attach boxed title to top left,
   % Ensures proper line breaking in longer titles
   minipage boxed title,
% (boxed title style requires an overlay)
boxed title style={empty,size=minimal,toprule=0pt,top=4pt,left=3mm,overlay={}},
coltitle=colexam,fonttitle=\bfseries,
before=\par\medskip\noindent,parbox=false,boxsep=0pt,left=3mm,right=0mm,top=2pt,breakable,pad at break=0mm,
   before upper=\csname @totalleftmargin\endcsname0pt, % Use instead of parbox=true. This ensures parskip is inherited by box.
% Handles box when it exists on one page only
overlay unbroken={\draw[colexam,line width=.5pt] ([xshift=-0pt]title.north west) -- ([xshift=-0pt]frame.south west); },
% Handles multipage box: first page
overlay first={\draw[colexam,line width=.5pt] ([xshift=-0pt]title.north west) -- ([xshift=-0pt]frame.south west); },
% Handles multipage box: middle page
overlay middle={\draw[colexam,line width=.5pt] ([xshift=-0pt]frame.north west) -- ([xshift=-0pt]frame.south west); },
% Handles multipage box: last page
overlay last={\draw[colexam,line width=.5pt] ([xshift=-0pt]frame.north west) -- ([xshift=-0pt]frame.south west); }%
}{ex}



\def\prooftext{Proof} % If English
\NewDocumentEnvironment{Proof}{ O{} }
{
    \colorlet{colexam}{black!55!black} % Global example color
    \newtcolorbox[number within=section]{testproofbox}{%
        empty,% Empty previously set parameters
        title={\emph{\prooftext} \thetcbcounter: #1},
        attach boxed title to top left,
        minipage boxed title,
        boxed title style={empty,size=minimal,toprule=0pt,top=4pt,left=3mm,overlay={}},
        coltitle=colexam,
        fonttitle=\bfseries,
        before=\par\medskip\noindent,
        parbox=false,
        boxsep=0pt,
        left=3mm,
        right=0mm,
        top=2pt,
        breakable,
        pad at break=0mm,
        before upper=\csname @totalleftmargin\endcsname0pt,
        % Removed height constraints
        % Added flexible width
        width=\dimexpr\textwidth-3mm\relax,
        % Modified overlays
        overlay unbroken={\draw[colexam,line width=.5pt] ([xshift=-0pt]title.north west) -- ([xshift=-0pt]frame.south west); },
        overlay first={\draw[colexam,line width=.5pt] ([xshift=-0pt]title.north west) -- ([xshift=-0pt]frame.south west); },
        overlay middle={\draw[colexam,line width=.5pt] ([xshift=-0pt]frame.north west) -- ([xshift=-0pt]frame.south west); },
        overlay last={\draw[colexam,line width=.5pt] ([xshift=-0pt]frame.north west) -- ([xshift=-0pt]frame.south west); },%
    }
    \begin{testproofbox}
}
{\end{testproofbox}\endlist}


%  ╔══════════════════════════════════════════════════════════╗
%  ║                    for Dynkin diagrams                   ║
%  ╚══════════════════════════════════════════════════════════╝

\def\row#1/#2!{#1_{\IfStrEq{#2}{}{n}{#2}} & \dynkin{#1}{#2}\\}
\newcommand{\tble}[1]{
   \renewcommand*\do[1]{\row##1!}
   \[
      \begin{array}{ll}\docsvlist{#1}\end{array}
   \]
}

%  ╔══════════════════════════════════════════════════════════╗
%  ║                         links                            ║
%  ╚══════════════════════════════════════════════════════════╝

\definecolor{LinkColour}{HTML}{A64A3B} % favourite
% \definecolor{LinkColour}{HTML}{8C3D32} % 2nd place
\hypersetup{
  colorlinks,
  citecolor=black,
  filecolor=magenta,
  linkcolor=LinkColour,
  urlcolor=blue,
}


%  ╔══════════════════════════════════════════════════════════╗
%  ║                          misc                            ║
%  ╚══════════════════════════════════════════════════════════╝

\stackMath %to be able to stack text in math mode
\makeindex

%import options: all, bibliography, book-formatting, booksSetting, customEnvironments, dynkin, hw-formatting, language-setup, newcommands, packages, preamble, proofs, settings, tcolorbox, test.svg,
\begin{document}
\maketitle

Intuition here

\section{Build up}


Let us start with some useful results from representation theory:

\begin{thm}{Clifford's Theorem}{cliffordThm}
	$G$ be a group, $N\nsubg G$ a normal subgroup, $K$ a field, $\pi: G \to \GL_n(K)$ be an irreducible
	representation. Then the restriction $\pi|_N$ breaks up into a direct sum of irreducible representations
	of $N$ of equal dimensions (i.e. it is semisimple and each representation is of the same dimension). These
	irreducible representations of N lie in one orbit for the action of G by conjugation on the equivalence
	classes of irreducible representations of N (i.e. they are isotypic components with the same multiplicity).
\end{thm}

\begin{Proof}
	Let $W$ be an irreducible subrepresentation of $V|_H$. Since $H$ is normal in $G$, for $g \in G$, $H$ acts
	irreducibly on $gW$ by $(h, gw) \mapsto hgh^{-1}hw$. The subspace $\sum_{g \in G} gW$ is a nonzero
	subrepresentation of $V$. Since $V$ is irreducible, it is equal to $V$. Since a representation which is a sum
	of irreducible sub-representations is semisimple, $V|_H$ is semisimple. A similar argument follows for the 2nd
	assertion, completing the proof.
\end{Proof}

Thus, assume $H\nsubg G$ is a finite index normal subgroup, and let $\pi$ be an irreducible $K$-representation
of $H$. Let us show that it is a component of an $K$-representation $\Pi$ of $G$ whose restriction to $H$
contains $\pi$. First, induct so that we get $\Ind_H^G\pi$. This is finite length still as $[G:H] < \infty$,
and contains an irreducible representation, say $\Pi$. Then there is a non-trivial map $\Pi \to \Ind_H^G
\pi$. By Frobenius:
\[
  \Hom_G(\Pi, \Ind_H^G, \pi) \cong \Hom_H(\Pi|_H, \pi)
\]
Hence, as $\pi$ is irreducible and the restriction is semisimple, $\pi$ must be in the direct summands of
$\Pi|_H$. Next, if $\chi$ is an $K$-character of $G$ that is trivial on $H$, the restriction of $\Pi\otimes
\chi$ to $H$ contains $\pi$, for
\[
  (\Pi \otimes \chi)|_H \cong \Pi|_H
\]
If there is another irreducible $K$-representation $\Pi'$ that contains $\pi$, it is a twist of $\Pi$:

\begin{lem}{lemma 1}{lem1}
    Let $K$ algebraically closed and \(G/H\) abelian. Then Any irreducible \(K\)-representation \(\Pi'\) of
    \(G\) containing \(\pi\) is isomorphic to \(\Pi \otimes \chi\) for some \(K\)-character \(\chi\) of \(G\) trivial on \(H\).
\end{lem}

\begin{Proof}
    Certainly \(\mathrm{Hom}_H(\Pi'|_H, \Pi|_H) \neq 0\). The right adjoint
    of the restriction from \(G\) to \(H\) is the induction \(\mathrm{Ind}_H^G\) from \(H\) to \(G\),
    therefore \(\Pi'\) is isomorphic to an irreducible subrepresentation of \(\mathrm{Ind}_H^G(\Pi|_H)\). We
    have \(\mathrm{Ind}_H^G(\Pi|_H) \simeq (\mathrm{Ind}_H^G 1) \otimes \Pi\) because \(G/H\) is finite, and
    the irreducible sub quotients of \(\mathrm{Ind}_H^G 1\) are the characters \(\chi\) of \(G\) trivial on
    \(H\) because \(K\) is algebraically closed. Therefore, there exists \(\chi\) such that \(\Pi' \simeq
    \Pi\otimes \chi\), completing the proof.
\end{Proof}

Suppose $H$ is a closed subgroup of a locally profinite group $G$ and $V$ a $K$-representation of $G$.
If $[G:H]<\infty$, $H$ is open, and conversely if $H$ is open and cocompact the index if finite. If $V$ is
smooth, $V|_H$ is smooth and if $H$ is open and $V|_H$ is smooth the $V$ is smooth.

Let $V$ be irreducible and $V|_H$ have finite length. Let $W$ be an irreducible component of $V|_H$, $\pi$ an
isomorphism class, and $G_\pi$ the $G$-stabilizer of $\pi$. Let $V_\pi$ be the $\pi$-isotypic component of
$V|_H$. Then the $G$-stabilizer of $V_\pi$ is $G_\pi$. The $G$-stabilizer of $W$ is open in $G$ as it contains
the nonzero $G$-stabilizer of $v\in W$ and $V$ is smooth, and furthermore it is contained in $G_\pi$. Both of
these groups have finite index in $G$ and by Clifford's theory $V = \Ind_{G_\pi}^G(V_\pi)$. The representation
$G_\pi$ on $V_\pi$ is irreducible and the length of $V|_H$ is:
\[
  \len(V|_H) = [G:G_\pi]\len(V_\pi|_H)
\]
\begin{lem}{lemma 2}{lem2}
  Assume that \(G/H\) is abelian. Then:
  \begin{enumerate}
  \item \(G_{\pi}\) is normal in \(G\) and does not depend on the choice of \(\pi\) in \(V|_{H}\). The smooth
    \(K\)-characters of \(G\) trivial on \(G_{\pi}\) are in \(X_V\).
  \item Assume \(K\) algebraically closed.
  \begin{enumerate}
  \item Any irreducible subquotient of the smooth induction \(\mathrm{Ind}_H^G 1\) is a smooth \(K\)-character \(\chi\) of \(G\) trivial on \(H\).
  \item Any irreducible \(K\)-representation of \(G\) containing \(\pi\) is a twist \(V \otimes \chi\) of
    \(V\) by some smooth \(K\)-character \(\chi\) of \(G\) trivial on \(H\).
  \item When \(V|_H\) has multiplicity 1, then \(W = V_{\pi}\), for a smooth \(K\)-character \(\chi\) of \(G\) trivial on \(H\), \(V \otimes \chi \simeq V\) if and only if \(\chi\) is trivial on \(G_{\pi}\), and \(G_{\pi}\) is the largest subgroup \(I\) of \(G\) containing \(H\) such that \(\mathrm{lg}(V|_I) = \mathrm{lg}(V|_H)\).
  \item When \(K\) is algebraically closed and \(V|_H\) has multiplicity 1, then
  \[
    |X_V| = \begin{cases} [G:G_{\pi}] & \text{if } \mathrm{char}_K = 0, \\ [G:G_{\pi,\ell}] & \text{if
    } \mathrm{char}_K = \ell > 0, \end{cases}
  \]
  where \(G_{\pi,\ell}\) is the smallest subgroup of \(G\) containing \(G_{\pi}\) such that \([G:G_{\pi,\ell}]\) is relatively prime to \(\ell\).
  \end{enumerate}
  \end{enumerate}
\end{lem}

\begin{Proof}
\begin{enumerate}
\item The isotypic components of \(\Pi|_H\) are \(G\)-conjugate, their \(G\)-stabilizers are \(G\)-conjugate and contain \(H\) hence they are equal because \(G/H\) is abelian. Since \(V \otimes \chi \simeq \mathrm{Ind}_{G_{\pi}}^G(\chi|_{G_{\pi}} \otimes V_{\pi})\) for any smooth \(K\)-character \(\chi\) of \(G\), the smooth \(K\)-characters of \(G\) trivial on \(G(\pi)\) are in \(X_V\).

\item
\begin{enumerate}
\item[(a)] For any closed subgroup \(Q\) of \(G\) and a smooth \(K\)-representation \(X\) of \(Q\), the representation \(\mathrm{Ind}_Q^G X\) is the space of functions \(f: G \to X\) with the property \(f(qgk) = qf(g)\) for \(q \in Q\), \(g \in G\), \(k \in K_f\) for some open subgroup \(K_f\) of \(G\), with the action of \(G\) by right translation, and where \(\mathrm{ind}_Q^G 1\) is the subrepresentation on the subspace of functions of compact support modulo \(Q\). When \(G/Q\) is compact, \(\mathrm{Ind}_Q^G X = \mathrm{ind}_Q^G X\).

Let \(V \supset U\) be \(G\)-stable subspaces with \(V/U\) irreducible. We can suppose \(V\) generated by an element \(f\) (indeed \(V'/U' \simeq V/U\) for the \(G\)-stable space \(V'\) generated by \(f \in V \setminus U\) and the kernel \(U'\) of the map \(V' \to V/U\)). There is an open subgroup \(K\) of \(G\) which fixes \(f\). We have \(U \subset V \subset \mathrm{ind}_K^G 1\) and one is reduced to the case where \(G/H\) is finite.

\item[(b)] The proof of Lemma \ref{lm:lem1} remains valid with the smooth induction \(\mathrm{Ind}_H^G\), which is the smooth compact induction \(\mathrm{ind}_H^G 1\), because \(G/H\) is compact, so that \(\mathrm{ind}_H^G(\Pi|_H) = \Pi \otimes \mathrm{ind}_H^G 1\).
\end{enumerate}

\item Any smooth character \(\chi\) of \(G\) trivial on \(H\) with \(\mathrm{ind}_{G_{\pi}}^G(V_{\pi}) \simeq \mathrm{ind}_{G_{\pi}}^G(V_{\pi} \otimes \chi|_{G_{\pi}})\) is trivial on \(G_{\pi}\). Indeed, restricting to \(G_{\pi}\) we see that \(V_{\pi} \otimes \chi|_{G_{\pi}}\) is conjugate to \(V_{\pi}\) by some \(g \in G\). Restricting to \(H\) gives that \(\pi \simeq \pi^g\), so \(g \in G_{\pi}\), hence \(V_{\pi} \otimes \chi|_{G_{\pi}} \simeq V_{\pi}\). As \(\mathrm{Ker}(\chi)\) is open in \(G\) and \(G/H\) is compact, \(J = \mathrm{Ker}(\chi) \cap G_{\pi}\) has finite index in \(G_{\pi}\). If \(\chi\) is not trivial on \(G_{\pi}\) then the action of \(J\) on \(V_{\pi}\) is reducible. Indeed, \(\mathrm{ind}_{J}^{G_{\pi}}(1)\) contains subrepresentations \(1\) and \(\chi|_{G_{\pi}}\), and by Frobenius reciprocity \(\mathrm{End}_J(V_{\pi}|_J)\) is equal to \(\mathrm{Hom}_{G_{\pi}}(V_{\pi}, \mathrm{ind}_J^{G_{\pi}}(V_{\pi}|_J)) = \mathrm{Hom}_{G_{\pi}}(V_{\pi}, V_{\pi} \otimes \mathrm{ind}_J^{G_{\pi}}(1))\).

Hence \(\dim(\mathrm{End}_{J}(V_{\pi}|_J)) \ge 2\) and \(V_{\pi}|_J\) is reducible. By the hypothesis of multiplicity 1, \(V_{\pi}|_H\) is irreducible, hence \(V_{\pi}|_J\) is irreducible as \(H \subset J\). So \(\chi\) is trivial on \(G_{\pi}\).

The group \(G_{\pi}\) is a subgroup \(I\) of \(G\) containing \(H\) with \(\mathrm{len}(V|_I) = \mathrm{len}(V|_H)\). If \(I\) has this property, the restriction to \(H\) of any irreducible component on \(V|_I\) is irreducible, hence \(I\) is contained in \(G_{\pi}\).

\item follows from (3).
\end{enumerate}
\end{Proof}

\begin{titledBox}{Restriction and Multiplicity}{resToMult}
Assume that \(V|_H\) has multiplicity 1. The \(G\)-stabilizer of any irreducible component of \(V\) is \(G_{\pi}\). Denote \(G_{\pi} = G_V\). Let \(I\) be a subgroup of \(G\) containing \(H\). The number of orbits of \(I\) in the irreducible components of \(V|_{G_V}\) is \(\mathrm{len}(V|_I)\). This number is the same for \(I\) and \(IG_V\), hence \(\mathrm{len}(V|_I) = \mathrm{len}(V|_{IG_V})\). We deduce that \(G_V \subset I\) if \(V|_I\) is reducible and \(|G/I|\) is a prime number.

\end{titledBox}

\begin{lem}{lemma 3}{lem3}
Let \(\theta\) be a smooth \(K\)-representation of a closed subgroup \(U \subg H\). Then if the functor \(\mathrm{Hom}_U(-, \theta)\) is exact on smooth \(K\)-representations of \(H\) and \(\dim \mathrm{Hom}_U(V, \theta) = 1\), then \(V|_H\) has multiplicity 1.
\end{lem}
\begin{Proof}
We denote by \(m_V(\pi)\) the multiplicity of any irreducible smooth \(K\)-representation \(\pi\) of \(H\) in
\(V|_H\). By exactness,
\[
	\sum_{\pi} m_V(\pi) \dim \mathrm{Hom}_U(\pi, \theta) = \dim \mathrm{Hom}_U(V, \theta) = 1.
\]
There is a single \(\pi\) with \(m_V(\pi) = \dim \mathrm{Hom}_U(V, \theta) = 1\).
\end{Proof}

\section{p-adic Reductive Group}

Need theorem 3.2 and corollary 3.5

Let $G$ be a $p$-adic (algebraic) reductive group, that is
the group of $F$-rational points of a reductive connected $F$-group $\underline{G}$, so $G=
\underline{G}(F)$.

Insure to incorporate this: \url{https://en.wikipedia.org/wiki/Reductive_group}




\begin{thm}{}{3.2}
  Let \(f: \underline{H} \to \underline{G}\) be an \(F\)-morphism of reductive connected \(F\)-groups such
  that \(f(H)\) is a normal subgroup of \(G\) of compact quotient \(G/f(H)\). Then:

  \begin{enumerate}
    \item  The restriction to \(f(H)\) of any irreducible admissible \(K\)-representation of \(G\) is semisimple of finite length.
    \item Any irreducible admissible \(K\)-representation of \(f(H)\) is contained in some irreducible
      admissible \(K\)-representation of \(G\) restricted to \(f(H)\), and extends to an irreducible
      admissible representation of some open subgroup of \(G\) of finite index.
  \end{enumerate}
\end{thm}

\begin{Proof}
  \(G\) satisfies (2-1) and \(f(H)\) satisfies the property (2-2) because \(H\) does. Apply the results of Section 2.2.
\end{Proof}

\begin{titledBox}{Remark 3.4}{rmk3.4}
The condition that f is central in \textcolor{blue}{Proposition 3.3} is necessary. Indeed, assume \(\mathrm{char}_F = 2\) and \(f: \underline{\mathrm{GL}}_2 \to \underline{\mathrm{SL}}_2\), \(f(g) = \varphi(g) / \det g\) where \(\varphi(x) = x^2\) for \(x \in F\) is the Frobenius.\textsuperscript{4} The \(F\)-morphism \(f\) is surjective but not central. Let \(G = \mathrm{GL}_2(F)\), \(G' = \mathrm{SL}_2(F)\), \(T'\) the diagonal torus of \(G'\) and \(U\) the group of unipotent upper triangular matrices in \(G'\). Then \(f(G) = T'\varphi(G')\) is closed but not normal and not cocompact in \(G'\) (since \(\varphi(U) = U \cap T'\varphi(G')\) and \(U/\varphi(U)\) homeomorphic to \(F/F^2\) is not compact).
\end{titledBox}

\begin{cor}{}{3.5}
Assume \(K\) algebraically closed. Let \(f: \underline{H} \to \underline{G}\) be an \(F\)-morphism of connected reductive \(F\)-groups which induces a central \(F\)-isogeny \(\underline{H}^{\mathrm{der}} \to \underline{G}^{\mathrm{der}}\) between the derived groups. Then the conclusions of \textcolor{blue}{Theorem 3.2} apply to \(f(H)\).
\end{cor}
\begin{Proof}
The \(F\)-isogeny \(\underline{H}^{\mathrm{der}} \to \underline{G}^{\mathrm{der}}\) is surjective with finite kernel contained in the centre of \(\underline{H}^{\mathrm{der}}\) \textcolor{red}{[Springer 1998, §12.2.6]}. If \(\underline{Z}\) is the connected centre of \(\underline{G}\), the natural map \(\underline{Z} \times \underline{G}^{\mathrm{der}} \to \underline{G}\) is surjective \textcolor{red}{[Springer 1998, Corollary 8.1.6]}. Hence the obvious map \(\underline{Z} \times \underline{H} \to \underline{G}\) is surjective and central. \textcolor{blue}{Proposition 3.3} applies to \(Zf(H)\). But \(K\) being algebraically closed, \(Z\) acts by a character in any irreducible smooth \(K\)-representations of \(G\), and we get the corollary.
\end{Proof}

\begin{titledBox}{Remark 3.6}{rmk3.6}
There is a more elementary proof that the restriction to \(f(H)\) of any irreducible admissible \(K\)-representation of \(G\) is semisimple of finite length in \textcolor{red}{[Silberger 1979]}.
\end{titledBox}




\section{Restriction of $GL_2(F)$ to $SL_2(F)$}

Let $F$ be a locally compact nonarchimedean field with residue characteristic (ex. $F/\Q_p$), and $L
= \overline{L}$ an algebraically closed field of characteristic $\ne p$. This paper will be about
$K$-representations of of $\SL_2(F)$. The reader can assume $K = \C$ to link this to the intuition complex
representations like in~\cite{bushnellLocalLanglandsConjecture2006}. The goal is to show how the Local Langlands
correspondence goes from being a bijection to a surjection, with fibers of size $1$, $2$, or $4$; these fibers
are called \emph{L-packets} or \emph{L-indistinguishable's}.

Let $G = \GL_2(F)$, $G' =\SL_2(F)$, $Z$ the center, $G$ the torus, $B$ the Borel subgroup, $U$ the unipotent
subgroup. If $X \sse G$ is a subset, $X' = X\cap G'$ and $x = (x_{ij})$ represents a matrix in $G$

Take $H= ZG'$. It is a closed normal compact group where $G/H$ is isomorphic to $F^*/(F^*)^2$. It is
normal as $Z$ is obviously normal and $G' = \ker\det$ where $\im\det = F^*$, and $gZG'g^{-1} = ZG'$. It is
closed since $H= \det^{-1}\left((F^*)^2\right)$. It is cocompact since:
\[
  \frac{F^*}{\left(F^*\right)^2} \cong \frac{\Z}{2\Z} \times \frac{\CO_F^*}{(\CO_F^*)^2}
\]
and $\CO_F^*$ is compact and so is the quotient by a closed subgroup The $\F_2$-dimension of $G/H$ is that of
$F^*/(F^*)^2$ which is:
\[
\dim_{\mathbb{F}_2} F^* / (F^*)^2 =
\begin{cases}
2 + e & \text{if } \text{char}_F \neq 2, \\
\infty & \text{if } \text{char}_F = 2,
\end{cases}
\quad \text{where } 2\CO_F = \FP_F^e.
\]
Hence if $\Char_F \ne 2$, $ZG'$ is open. Next, let $\overline{F}^\sep$ represent the separable algebraic
closure, $\CW_F$ and $G_F$ the Weil group and galois group of $\overline{F}^\sep/F$ respectively. Let $\Psi$
be a smooth additive $K$-character on of $F$, trivial on $\CO_F$ but not on $\fp_F^{-1}$ (i.e. a level
1 character, but not on $\fp_F^{-1}$ (i.e. A level 1 character).


\index{Whittaker Space}
\begin{defn}{Whittaker Space}{whittakerSp}
  The space of smooth $K$-characters of $U$ have the form
  \[
    \theta(u) = \psi \circ \tr(Y(u-1)) = \Psi(Y_{2,1}, u_{1,2}) \qquad u\in U
  \]
  for a lower triangular nilpotent matrix $Y \in M_2(F)$. The case $Y=0$ gives the trivial character of $U$,
  the cases with $Y\ne 0$ give the \emph{nondegenerate characters} of $U$.
\end{defn}

When $Y_{2,1} = 1$, then \textcolor{red}{here}, so denote $\theta_Y = \theta$.


\textcolor{red}{mores words here before defining}

\index{L-packet}
\begin{defn}{L-packet}{l-packet}
Let \(\Pi\) be an irreducible smooth \(K\)-representation of \(G\). Then the set \(L(\Pi)\) of isomorphism classes
of irreducible components of the restriction \(\Pi|_{G'}\) is called an \emph{L-packet}.
\end{defn}

We classify the irreducible smooth \(K\)-representations of \(G'\) by these L-packets.

\begin{prop}{Properties of L-packets: I}{propofL-packetsI}
  \begin{enumerate}
    \item  Any irreducible smooth \(K\)-representation of \(G'\) belongs to a unique \(L\)-packet.

    \item For two irreducible smooth \(K\)-representations \(\Pi_1, \Pi_2\) of \(G\),
  \[
    L(\Pi_1) = L(\Pi_2) \iff \Pi_1 = (\chi \circ \det) \otimes \Pi_2
  \]
  for some \(K\)-character \(\chi \circ \det\) of \(G\).

  \item, The trivial character of \(G'\) is the unique finite-dimensional irreducible smooth \(K\)-representation of \(G'\), it is degenerate and forms an \(L\)-packet \(L(1) = L(\chi \circ \det)\) for any smooth \(K\)-character \(\chi\) of \(F^*\).

\item If \(\Pi\) is an irreducible smooth \(K\)-representation of \(G\) the restriction of \(\Pi\) to \(G'\)
  is semisimple of finite length and multiplicity 1.

  \end{enumerate}
\end{prop}

\begin{Proof}
here
\end{Proof}

\textcolor{red}{words here before}

\begin{prop}{Properties of L-packets: II}{propofL-packetsII}
  \begin{enumerate}
  \item \(G_{\Pi}\) is the largest subgroup \(I\) of \(G\) such that \(\mathrm{len}(\Pi|_I) = \mathrm{len}(\Pi|_{G'})\).
  \item \(\Pi = \mathrm{ind}_{G_{\Pi}}^G V_{\pi}\) where \(V_{\pi}\) is the space of \(\pi\).
  \item \(L(\Pi)\) is a principal homogeneous space for \(G/G_{\Pi}\).
  \item \(|L(\Pi)|\) is a power of 2, and \(|L(\Pi)|\) divides \(2^{2+e}\) when \(\mathrm{char}_F \neq 2\).
  \end{enumerate}
\end{prop}

\begin{Proof}
here
\end{Proof}

If $p$ is odd, then $|F^*/(F^*)^2| = 4$, and hence:

\begin{prop}{Size of L-packets}{sizeofL-packets}
  If $p$ is odd, then the cardinality of an $L$-packet is either $1$, $2$, or $4$.
\end{prop}

\begin{Proof}
here
\end{Proof}

When $p =2$, this more difficult as \textcolor{red}{HERE}. This is proven via LLC.




















% \newpage
% \printbibliography

\end{document}
